\pdfvariable trailerid {[<C3662026090400000000000000000000><C3662026090400000000000000000000>]}
\documentclass[10pt]{article}
\usepackage[margin=0.68in]{geometry}
\usepackage{amsmath,amssymb,amsthm,mathtools,microtype}
\usepackage[hidelinks]{hyperref}
\usepackage{lmodern}
\setlength{\parindent}{1em}
\setlength{\parskip}{0.15em}
\emergencystretch=3em
\newtheorem{theorem}{Theorem}
\newtheorem{proposition}{Proposition}
\renewcommand{\qedsymbol}{\textsc{qed}}
\providecommand{\CRevisionRound}{2}
\newcommand{\C}{\mathbb C}

\title{A Krawtchouk XX Chain: Exact Propagation, Perfect Mirror Transfer,\\
and the Full Fermionic Phase}
\author{HCS-C366 source-local theorem package}
\date{September 4, 2026}

\begin{document}
\maketitle

\begin{abstract}
We identify the engineered open XX chain with the spin-
\(N/2\) representation of \(J_x\), obtaining its complete Krawtchouk spectrum
and propagator for every chain length.  An endpoint excitation follows an
exact binomial law and is mirrored perfectly at time \(\pi/\Omega\).
\ifnum\CRevisionRound>0
Fermionic second quantization then gives the propagator in every excitation
sector and exposes the wedge-reordering sign that a one-particle argument
cannot determine.
\fi
\ifnum\CRevisionRound>1
We close all size, coupling, field, and revival boundaries and distinguish a
sectorwise phase from a global Fock-space revival.
\fi
Finite exact ledgers audit the formulas; representation theory proves the
all-size statements.
\end{abstract}
\ifnum\CRevisionRound>1
\small
\fi

\section{Frozen chain and one-particle spectrum}
On the ordered basis \(\{|j\rangle:0\le j\le N\}\), let
\begin{equation}
h_N=\frac{\Omega}{2}\sum_{j=0}^{N-1}\sqrt{(j+1)(N-j)}
 (|j\rangle\langle j+1|+|j+1\rangle\langle j|),\qquad \Omega>0.
\end{equation}
Identifying \(|j\rangle\) with the \(J_z\)-weight \(j-N/2\), the standard
matrix element of \(J_x\) makes (1) exactly \(\Omega J_x\).

\begin{theorem}[Krawtchouk spectral resolution]
The simple eigenvalues are
\begin{equation}
E_r=\Omega(N/2-r),\qquad 0\le r\le N.
\end{equation}
If
\begin{equation}
K_r(j)=\sum_\ell(-1)^\ell\binom j\ell\binom{N-j}{r-\ell},
\end{equation}
then normalized eigenvectors have components
\begin{equation}
v_r(j)=\frac{\sqrt{\binom Nj}K_r(j)}{\sqrt{2^N\binom Nr}}.
\end{equation}
\end{theorem}

\begin{proof}
Spin rotation sends the \(J_z\) weights to the \(J_x\) weights, proving
(2).  Pascal identities in the tridiagonal recurrence prove
\(h_Nv_r=E_rv_r\).  Coefficient extraction from
\[
\sum_j\binom Nj(1-z)^j(1+z)^{N-j}(1-w)^j(1+w)^{N-j}
=2^N(1+zw)^N
\]
gives
\(\sum_j\binom NjK_r(j)K_s(j)=2^N\binom Nr\delta_{rs}\).
Thus the \(N+1\) displayed vectors are a complete orthonormal basis.
\end{proof}

\begin{theorem}[Exact endpoint propagation]
For \(U_1(t)=e^{-ith_N}\),
\begin{equation}
\langle k|U_1(t)|0\rangle=(-i)^k\sqrt{\binom Nk}
 \sin^k\!\frac{\Omega t}{2}\cos^{N-k}\!\frac{\Omega t}{2}.
\end{equation}
At \(t_*=\pi/\Omega\), \(U_1(t_*)=(-i)^NR\), where
\(R|j\rangle=|N-j\rangle\).
\end{theorem}

\begin{proof}
The spin-\(N/2\) module is the symmetric power of \(N\) spin halves.
Expanding
\(e^{-it\Omega\sigma_x/2}|0\rangle=\cos(\Omega t/2)|0\rangle-
i\sin(\Omega t/2)|1\rangle\) in normalized symmetric weights gives (5).
At angle \(\pi\), every factor flips with phase \(-i\).
\end{proof}
This identity is the \emph{single-particle propagator owner}.  In particular,
at \(t_*/2\) the endpoint-start probabilities are \(2^{-N}\binom Nk\).

\ifnum\CRevisionRound>0
\section{The many-body theorem}
Let \(H_N=d\Gamma(h_N)\), and order every fermionic basis wedge increasingly.

\begin{theorem}[Fock-space mirror and complete energies]
On the \(m\)-particle sector,
\begin{equation}
e^{-itH_N}=\bigwedge^m U_1(t).
\end{equation}
At \(t_*\),
\begin{equation}
|j_1<\cdots<j_m\rangle\mapsto
(-i)^{mN}(-1)^{m(m-1)/2}
|N-j_m<\cdots<N-j_1\rangle .
\end{equation}
For \(S\subseteq\{0,\ldots,N\}\), \(|S|=m\), the corresponding energy is
\begin{equation}
E_S=\Omega\left(mN/2-\sum_{r\in S}r\right),
\end{equation}
and all degeneracies are encoded by
\begin{equation}
\prod_{r=0}^N(1+yq^r)=
\sum_{m=0}^{N+1}y^mq^{m(m-1)/2}
\mathcal G_{N+1,m}(q).
\end{equation}
Here the Gaussian $q$-binomial polynomial is fixed by
\begin{equation}
\mathcal G_{n,m}(q)=\mathcal G_{n-1,m}(q)+
q^{n-m}\mathcal G_{n-1,m-1}(q),
\qquad
\mathcal G_{n,0}(q)=\mathcal G_{n,n}(q)=1,
\end{equation}
with out-of-range entries zero.
\end{theorem}

\begin{proof}
A number-conserving quadratic fermion Hamiltonian induces the exterior power
of its one-particle unitary.  At \(t_*\), each occupied orbital contributes
\((-i)^N\).  Reflection reverses all \(m\) wedge factors, and restoring
increasing order needs \(m(m-1)/2\) transpositions.  Exterior products of the
one-particle eigenbasis prove (8).

Let $A_{N,m}(q)$ be the coefficient of $y^m$ in the product in (9).
Splitting subsets according to whether they contain $N$ gives
$A_{N,m}=A_{N-1,m}+q^N A_{N-1,m-1}$.  The expression
$q^{m(m-1)/2}\mathcal G_{N+1,m}(q)$ satisfies the same recurrence:
the second term of (10) supplies $q^{N+1-m}$, and shifting the prefactor from
$m$ to $m-1$ supplies $q^{m-1}$.  The empty product is the common base case,
so induction proves (9) coefficient by coefficient.
\end{proof}
Equation (7) is the \emph{many-body phase owner}; deleting its reordering sign
changes the unitary in entire excitation sectors.
\fi

\ifnum\CRevisionRound>1
\section{Boundary and revival audit}
For \(N=0\), reflection is the one-site identity.  For \(\Omega=0\), the
Hamiltonian vanishes and distinct endpoints do not transfer.  Negative
\(\Omega\) reverses time and replaces \((-i)^N\) by \(i^N\) at positive
time \(\pi/|\Omega|\).  A uniform field \(B\widehat m\) contributes only
\(e^{-imBt}\) in sector \(m\).

For the field-shifted model the exact Fock operators are
\[
 U_B(2\pi/|\Omega|)=
 e^{-i(2\pi B/|\Omega|)\widehat m}(-1)^{N\widehat m},\qquad
 U_B(4\pi/|\Omega|)=e^{-i(4\pi B/|\Omega|)\widehat m}.
\]
The first is identity exactly when $2B/|\Omega|+N\in2\mathbb Z$; the second
is identity exactly when $2B/|\Omega|\in\mathbb Z$.  In particular, for the
unshifted model $B=0$, the $2\pi/|\Omega|$ propagator is identity for even
$N$ and fermion parity for odd $N$, while $4\pi/|\Omega|$ is always identity.
These conditions include the vacuum sector, so a sectorwise phase is not
silently promoted to a global Fock revival.  This is the
\emph{revival-boundary owner}.  No persistence of perfect transfer under a
generic coupling perturbation is asserted.

\ifnum\CRevisionRound>1
\small
\fi
\section{Evidence and Route-A boundary}
The executable receipt covers 66 spectral rows, all 65,534 fermionic subsets
through \(N=14\), complete energy histograms, 231 formal all-time endpoint
cells, and all 136 Gaussian-polynomial rows through order 15.
Independent code and SymPy reconstruct the formulas; finite checks do not
replace the proofs.

The strict tuple is
\[
(\mathrm{A0\_FAIL},\mathrm{A1\_WEAK},\mathrm{A2\_FAIL},
\mathrm{A3\_FAIL},\mathrm{A4\_NATURAL\_QUANTIZATION}).
\]
Natural quantization does not automatically furnish arithmetic ownership or
an analytic Euler-product mechanism.
Overall, Route A is rejected and Route B is false.  The finite source
Hamiltonian has no rational-prime owner, logarithmic-prime clock, target
Euler product, divisor, functional equation, zero match, or Hilbert--P\'olya
interpretation.  Scope is \texttt{NO\_BAD\_EULER\_OR\_ROOT\_NUMBER}.
\fi

\ifnum\CRevisionRound>1
\enlargethispage{3\baselineskip}
\fi
\section*{Source and claim boundary}
Christandl, Datta, Ekert, and Landahl introduced fixed-coupling perfect
transfer networks in \emph{Physical Review Letters} 92 (2004), 187902,
\href{https://doi.org/10.1103/PhysRevLett.92.187902}{doi:10.1103/PhysRevLett.92.187902}.
Albanese, Christandl, Datta, and Ekert treated mirror inversion in
\emph{Physical Review Letters} 93 (2004), 230502,
\href{https://doi.org/10.1103/PhysRevLett.93.230502}{doi:10.1103/PhysRevLett.93.230502}.
They establish lineage; the frozen formulas above are rederived, and no
literature-priority claim is made.
\ifnum\CRevisionRound>1
\normalsize
\fi

\end{document}
